\section{Related Work}
\textit{Overhead priors.} Airborne lidar and imagery at metre scale have served as route priors for ground vehicles~\cite{vandapel2006aerial}; drones can map ahead of the robot~\cite{delmerico2017active} or learn terrain properties from the air~\cite{fortin2025uav}, and planetary missions plan from orbital data~\cite{ono2016traversability}. We ask what is lost when the prior is a coarse map, the only free kind for most deserts.

\textit{Slip on sand.} Geometric costmaps score slope and roughness on a grid~\cite{papadakis2013,goldberg2002}. On loose sand slip rises steeply with slope~\cite{ishigami2007,heverly2013,shrivastava2020}; slip can be learned from appearance~\cite{angelova2007slip}, desert vehicles have planned on overhead and onboard maps~\cite{thrun2006stanley}, and planners propagate slip uncertainty into routes~\cite{inotsume2020robust}, but they assume that the map resolves the slope.

\textit{Dunes in elevation models.} Dune morphometry from DEMs and imagery is well developed~\cite{hugenholtz2012remote}, and deep learning maps crestlines from DEM derivatives~\cite{daynac2024workflow}. Global 30\,m products have been ranked against lidar~\cite{guth2021ranking}; TanDEM-X, the radar behind GLO-30, is weaker over sand~\cite{martone2016enhancing}, AW3D30 comes from optical stereo~\cite{tadono2014aw3d}, and DEM error in a dune sea grows above 15\,degrees~\cite{wang2025multisource}. Super resolution can sharpen coarse DEMs~\cite{chen2016demsr}. We know of no benchmark that scores coarse map hazard layers and planners for slip faces against lidar truth across resolutions.
